\documentclass[a4paper,12pt]{article}

\oddsidemargin=-4mm
\evensidemargin=-4mm
\textheight=255mm
\topmargin=-15mm
\textwidth=170mm
\headsep=10mm

\usepackage{amssymb,amsmath,amsfonts}
\usepackage[T2A]{fontenc}
\usepackage[english,russian]{babel}
\usepackage{graphicx}
\usepackage{epstopdf}

\let\chapter\section
\sloppy

\begin{document}

%%%%%%%%%%%%%%%%%%%%%%%%%%%%%%%%%%%%%%%%%%%%%%%%%%%%%%%%%%%%%%%%%%%%%%%%%%%%%%%%%%%%%%%%%%%
\noindent
УДК: 532.546

{\large
\begin{center}
\textbf{Моделирование выпадения парафина при закачке холодной воды}\\
И.С.~Лазарев$^1$, А.И.~Никифоров$^2$
\end{center}
} 

{\footnotesize
\noindent $^1$ Lazarevigor.2001@mail.ru; КФУ, ИММ, им. Н.И.~Лобачевского \\
\noindent $^2$ nikiforov@imm.knc.ru; ФИЦ КазНЦ РАН, ИММ
}

%\hrule

\begin{abstract}
This paper describes a recursive approach to the enumeration of some classes of combinatorial tasks. Most tasks are used in the
specific scope of teaching and learning informatics through olympiads and other competitions. Combinatorial problems can often
lead to interesting and beautiful dynamic programming tasks, because they both depend on recurrence relations: formulae that
solve a larger problem in terms of one or more smaller problems.
\end{abstract}

\noindent
\textit{Keywords:} programming contests, informatics olympiads, combinatorial tasks.


\vspace{1em}
\hrule
\vspace{2em}
%%%%%%%%%%%%%%%%%%%%%%%%%%%%%%%%%%%%%%%%%%%%%%%%%%%%%%%%%%%%%%%%%%%%%%%%%%%%%%%%%%%%%%%%%%%





%%%%%%%%%%%%%%%%%%%%%%%%%%%%%%%%%%%%%%%%%%%%%%%%%%%%%%%%%%%%%%%%%%%%%%%%%%%%%%%%%%%%%%%%%%%

\bibliographystyle{gost2008}      % mathematics and physical sciences
\bibliography{LazarevIS}   % name your BibTeX data base

%%%%%%%%%%%%%%%%%%%%%%%%%%%%%%%%%%%%%%%%%%%%%%%%%%%%%%%%%%%%%%%%%%%%%%%%%%%%%%%%%%%%%%%%%%%

\begin{otherlanguage}{english}

{\large
\begin{center}
\textbf{Article's Title}\\
I.I.~Ivanov, P.P.~Petrov
\end{center}
}

\begin{abstract}
This paper describes a recursive approach to the enumeration of some classes of combinatorial tasks. Most tasks are used in the
specific scope of teaching and learning informatics through olympiads and other competitions. Combinatorial problems can often
lead to interesting and beautiful dynamic programming tasks, because they both depend on recurrence relations: formulae that
solve a larger problem in terms of one or more smaller problems.
\end{abstract}

\textit{Keywords:} programming contests, informatics olympiads, combinatorial tasks.

\end{otherlanguage}

%%%%%%%%%%%%%%%%%%%%%%%%%%%%%%%%%%%%%%%%%%%%%%%%%%%%%%%%%%%%%%%%%%%%%%%%%%%%%%%%%%%%%%%%%%%

\end{document}
